\QHeading{第1题　计算填表。}
\par\noindent


\DataTable{5}{商品 & 魔方 & 订书机 & 书包 & 水杯\\\hline
单价/元 & \(5\) & {\color{answerred}24} & \(45\) & \(18\)\\\hline
数量/个 & \(12\) & \(8\) & {\color{answerred}8} & \(20\)\\\hline
总价/元 & {\color{answerred}60} & \(192\) & \(360\) & {\color{answerred}360}\\\hline}
\par

\par\noindent

\Answer{
\QHeading{计算过程：计算填表}
\begin{center}
\begin{adjustbox}{max width=\linewidth}\begin{tabular}{c|rrrr}
\hline
商品&魔方&订书机&书包&水杯\\\hline
单价/元&5&24&45&18\\
数量/个&12&8&8&20\\
总价/元&60&192&360&360\\\hline
\end{tabular}\end{adjustbox}
\end{center}
根据“单价\(\times\)数量\(=\)总价”及其逆运算填表：
\begin{enumerate}
\item 魔方总价：\(\begin{aligned}[t]&\phantom{={}}5\times12\\&=5\times(10+2)\\&=5\times10+5\times2\\&=50+10\\&=60\text{（元）}\text{。}\end{aligned}\)
\item 订书机单价：\(\begin{aligned}[t]&\phantom{={}}192\div8\\&=(160+32)\div8\\&=160\div8+32\div8\\&=20+4\\&=24\text{（元）}\text{。}\end{aligned}\)
\item 书包数量：\(360\div45=8\)个，因为\(45\times8=360\)。
\item 水杯总价：\(\begin{aligned}[t]&\phantom{={}}18\times20\\&=18\times(2\times10)\\&=(18\times2)\times10\\&=36\times10\\&=360\text{（元）}\text{。}\end{aligned}\)
\end{enumerate}

}

\QHeading{第2题　先写出数量关系式，再列式解答。}

\Problem{①一个文具盒售价16元，买9个文具盒一共需要多少钱？}
\par


\MethodTitle{方法一}
\Answer{单价$\times$数量$=$总价。\par \Formula{\begin{aligned}[t]&\phantom{={}}16\times9\\&=16\times(10-1)\\&=16\times10-16\times1\\&=160-16\\&=144\text{（元）}\end{aligned}}\par 答：一共需要144元。}

\MethodTitle{方法二：拆成5个和4个}
\Answer{把9个文具盒分成5个和4个，分别算钱再合起来。}
\Answer{\Formula{\begin{aligned}[t]&\phantom{={}}16\times 9\\
&=16\times (5+4)\\
&=16\times 5+16\times 4\\
&=80+64\\
&=144\text{（元）}\end{aligned}}}
\Answer{答：一共需要144元。}
\Problem{②买7本相同的故事书一共花了91元，每本故事书多少钱？}

\MethodTitle{方法一}
\Answer{总价$\div$数量$=$单价。\par \Formula{\begin{aligned}[t]&\phantom{={}}91\div7\\&=13\text{（元）}\end{aligned}}\par 答：每本故事书13元。}

\MethodTitle{方法二：把总价拆成容易平均分的两部分}
\Answer{91元分成70元和21元，每部分都平均分给7本书。}
\Answer{\Formula{\begin{aligned}[t]&\phantom{={}}91\div 7\\
&=(70+21)\div 7\\
&=70\div 7+21\div 7\\
&=10+3\\
&=13\text{（元）}\end{aligned}}}
\Answer{答：每本故事书13元。}
\par\noindent

\QHeading{第3题　解决问题。}

\Problem{①新华文具店搞促销，每盒彩笔22元。小画家社团要买13盒彩笔，一共需要准备多少钱？}
\par


\MethodTitle{方法一}
\Answer{\Formula{\begin{aligned}[t]&\phantom{={}}22\times13\\&=22\times(10+3)\\&=22\times10+22\times3\\&=220+66\\&=286\text{（元）}\end{aligned}}\par 答：一共需要准备286元。}

\MethodTitle{方法二：先按每盒20元，再补每盒2元}
\Answer{先算13盒各20元的费用，再补每盒2元。}
\Answer{\Formula{\begin{aligned}[t]&\phantom{={}}22\times 13\\
&=(20+2)\times 13\\
&=20\times 13+2\times 13\\
&=260+26\\
&=286\text{（元）}\end{aligned}}}
\Answer{答：一共需要准备286元。}
\Problem{②学校采购绿植，花432元买了18盆同样的绿植，每盆绿植的单价是多少元？}

\MethodTitle{方法一}
\Answer{\Formula{\begin{aligned}[t]&\phantom{={}}432\div18\\&=(360+72)\div18\\&=360\div18+72\div18\\&=20+4\\&=24\text{（元）}\end{aligned}}\par 答：每盆绿植的单价是24元。}

\MethodTitle{方法二：把18盆平均分成两组}
\Answer{先算9盆的费用，再求每盆的费用。}
\Answer{\Formula{\begin{aligned}[t]&\phantom{={}}432\div 18\\
&=432\div (2\times 9)\\
&=432\div 2\div 9\\
&=216\div 9\\
&=24\text{（元）}\end{aligned}}}
\Answer{答：每盆绿植的单价是24元。}

